\documentclass[11pt,a4paper]{article}

%% PACHETE MATEMATICE
\usepackage{amsmath,amssymb,amsfonts}
\usepackage{amsthm}
\usepackage{mathpazo}
\usepackage{eulervm}
\usepackage{enumerate}
\usepackage{color}
\usepackage{graphicx}
\usepackage[all]{xy}
\usepackage{xcolor}
\usepackage{framed}
\usepackage[unicode]{hyperref}
\setcounter{MaxMatrixCols}{20}
\usepackage{multicol}
%\usepackage[romanian]{babel}


%% ASPECT
\linespread{1.05}
\usepackage[T1]{fontenc}
\usepackage[utf8x]{inputenc}
\usepackage[margin=2cm]{geometry}
% \usepackage[color]{showkeys}
\usepackage{anyfontsize}
\usepackage{titlesec}
\titleformat*{\section}{\large\bfseries}

\usepackage{fancyhdr}
\pagestyle{fancy}
\rhead{\small \color{gray} Notițe de Adrian Manea}
\lhead{\small \color{gray} M1-ACS, 2017---2018, M. Olteanu}


\usepackage[autostyle = false, style = german]{csquotes}
\usepackage[romanian]{babel}
\newcommand\qq{\enquote}

%% COMENZILE MELE
\renewcommand*{\proofname}{Dem.:}
\newcommand\bb{\mathbb}
\newcommand\dr{\mathrm}
\newcommand\kal{\mathcal}
\newcommand\fk{\mathfrak}
\newcommand\op{\oplus}
\newcommand\ot{\otimes}
\newcommand\ds{\displaystyle}
\newcommand\id{\indent}
\newcommand\nid{\noindent}
\newcommand\seq{\subseteq}
\newcommand\sneq{\subsetneq}
\newcommand\speq{\supseteq}
\newcommand\spneq{\supsetneq}
\newcommand\rar{\rightarrow}
\newcommand\Rar{\Rightarrow}
\newcommand\xrar{\xrightarrow}
\newcommand\lar{\leftarrow}
\newcommand\Lar{\Leftarrow}
\newcommand\xlar{\xleftarrow}
\newcommand\lrar{\leftrightarrow}
\newcommand\Lrar{\Leftrightarrow}
\newcommand\ideal{\trianglelefteq}
\newcommand\ai{\text{ a.\^{i}. }}
\newcommand\sk{(\textit{Schi\c{t}\u{a}}) }
\newcommand\rhup{\rightharpoonup}
\newcommand\rhdn{\rightharpoondown}
\newcommand\lhup{\leftharpoonup}
\newcommand\lhdn{\leftharpoondown}
\newcommand\vid{\emptyset}
\newcommand\wh{\widehat}
\newcommand\ol{\overline}
\newcommand\NN{\mathbb{N}}
\newcommand\RR{\mathbb{R}}
\newcommand\ZZ{\mathbb{Z}}
\newcommand\QQ{\mathbb{Q}}
\newcommand\CC{\mathbb{C}}

%% STILURILE MELE
\newtheoremstyle{dotless}{}{}{\itshape}{}{\bfseries}{:}{ }{}

\theoremstyle{dotless}
\newtheorem{theorem}{Teorem\u{a}}[section]
\newtheorem{proposition}{Propozi\c{t}ie}[section]
\newtheorem{lemma}{Lem\u{a}}[section]
\newtheorem{corollary}{Corolar}[section]
\newtheorem{exercise}{Exerci\c{t}iu}

\newtheoremstyle{dotlessdr}{}{}{}{}{\bfseries}{:}{ }{}
\theoremstyle{dotlessdr}
\newtheorem{example}{Exemplu}[section]
\newtheorem{definition}{Defini\c{t}ie}[section]
\newtheorem{remark}{Observa\c{t}ie}[section]




\begin{document}

\begin{center}
  {\large\textbf{Seminar 3 \\ Serii de numere reale}}
\end{center}

\textbf{Exerciții suplimentare}

1. Să se stabilească natura seriilor următoare:

\begin{enumerate}[(a)]
\item $\sum_n \frac{1}{\sqrt[n+1]{\ln (n + 1)}}$ (D, necesar [Stolz + $a^b = e^{b \ln a}$]);
\item $\sum_n \frac{1}{n} \sqrt[n]{(n + 1)(n + 2) \cdots (n + n)}$ (D, necesar + raport);
\item $\sum_n \sqrt{n^4 + 2n + 1} - n^2$ (D, comparație 3 cu $\sum \frac{1}{n}$);
\item $\sum_n \frac{1}{1 + \sqrt{2} + \sqrt[3]{3} + \dots + \sqrt[n]{n}}$ (D, comparație 3 cu $\sum \frac{1}{n}$ [Stolz]);
\item $\sum_n \frac{1}{\sqrt[n]{\ln n}}$ (D, comparație cu $\sum \frac{1}{\sqrt[n]{n}}$);
\item $\sum_n \frac{a^n}{\sqrt[n]{n!}}, a > 0$
  (comparație [$a = 1$] cu $\sum a^n$, raport pentru $a > 1$ și $\sum\frac{a^n}{n!}$, apoi comparație);
\item $\sum_n a^{\ln n}, a > 0$ ([$a = \frac{1}{e}$], Raabe);
\item $\sum_{n \geq 2} \frac{1}{n \ln n}$ (D, integral);
\item $\sum_n (-1)^n \frac{\log_a n}{n}, a > 1$ (C, Leibniz: $f(x) = \frac{\log_a x}{x}$ crescătoare pentru $x > e$);
\item $\sum_n \frac{(-1)^n \sqrt{n} + 1}{n}$ (D, spargem în două $\sum \frac{1}{n}$ D și restul C [Leibniz]);
\item $\sum_n \frac{1}{n^p \ln^q n}, p, q > 0$ ($p > 1$ C $\forall q > 0$ [comparație 1], $p = 1$ integral [C ddacă $q > 1$],
  $p < 1$ condensare $\frac{1}{n^q 2^{n(p-1)}} \ln^q 2$, D [raport]);
\item $\sum \frac{n!}{(a + 1)(a + 2) \cdots (a + n)}, a > 0$ (Raabe: $a > 1$, C, $a < 1$, D, $a = 1$, D [direct]);
\item $a + \sum_{n \geq 2} (2 - \sqrt{e})(2 - \sqrt[3]{e}) \cdots (2 - \sqrt[n]{e}) \cdot a^n, a > 0$
  (raport: $a < 1$, C, $a > 1$, D, comparație pentru $a = 1$ cu s.a.);
\item $\sum \frac{1}{n^a} \sin \frac{\pi}{n}, a \in \RR$ (comparație la limită, $\sum \frac{1}{n^{a + 1}}$);
\item $\sum \frac{a^n (n!)^2}{(2n)!}, a > 0$ (raport $\Rightarrow \frac{a}{4}$, $a = 4 \Rightarrow$ Raabe, D);
\item $\sum \frac{\cos n \cdot \cos \frac{1}{n}}{n}$ (Abel, $x_n = \frac{\cos n}{n}, y_n = \cos \frac{1}{n}$, C);
\item $\sum n^2 e^{-\sqrt{n}}$ (C, logaritmic);
\item $\sum n! \sin a \cdot \sin \frac{a}{2} \cdots \sin \frac{a}{n}, \ a \in (0, \pi)$. (raport, $a = 1$ Raabe + L'Hospital).
\end{enumerate}

\vspace{1cm}

2. Arătați că seria $\sum_n (-1)^n \frac{2 + (-1)^n}{n}$ este divergentă, dar
șirul termenilor converge la 0.

\vspace{1cm}

3. Studiați convergența seriei: $\sum_n \frac{\sin n \cdot \sin n^2}{\sqrt{n}}$
(Indicație: Abel pentru $\alpha_n = \frac{1}{\sqrt{n}}, v_n = \sin n \cdot \sin n^2$).

\vspace{1cm}

4. Studiați convergența absolută a seriilor:
\begin{enumerate}[(a)]
\item $\sum (-1)^n \sqrt[n]{n} \sin \frac{1}{n}$;
\item $\sum (-1)^{n + 1} \frac{2^n \sin^{2n} a}{n + 1}$, $a \in \RR$ (radical + Leibniz);
\item $\sum \frac{\sin n \cdot a}{n}, a \in \RR$ (comparație + Abel);
\item $\sum_n x_n$, unde $x_{2n - 1} = \frac{1}{\sqrt{n + 1} - 1}$ și $x_{2n} = -\frac{1}{\sqrt{n + 1} - 1}$;
\item $\sum_n x_n$, unde $x_{2n - 1} = \frac{1}{5n - 3}$ și $x_{2n} = -\frac{1}{5n - 3}$;
\end{enumerate}

\vspace{1cm}

5. Fie seria de termen general $x_n = \frac{(-1)^{n+1}}{\sqrt{n}}$. Arătați că ea este
semiconvergentă. Studiați seria obținută prin ridicarea ei la pătrat. Deduceți faptul că
este posibil ca produsul a două serii semiconvergente să fie o serie convergentă.

\vspace{1cm}

6. Considerați seriile:
\[
  S = 1 - \sum_n \Big( \frac{3}{2} \Big)^n, \quad %
  T = 1 + \sum_n \Big( \frac{3}{2} \Big)^{n - 1} \cdot \Big( 2^n + \frac{1}{2^{n + 1}} \Big).
\]
Arătați că seriile sînt divergente, dar produsul lor este o serie absolut convergentă.

\end{document}
